% Source: https://artofproblemsolving.com/community/c3h2021572_sum_to_infinity
% Stable identity: aops-topic-2021572-q01
By partial fractions,
\[
\dfrac{n}{(n+1)(n+2)}=\dfrac{2}{n+2}-\dfrac{1}{n+1}.
\]
This does not telescope to a finite limit, because the coefficients $2$ and $-1$ do not cancel. Instead compare with the harmonic series. For every $n\ge 1$ we have $(n+1)(n+2)=n^2+3n+2\le n^2+3n^2+2n^2=6n^2$, so
\[
\dfrac{n}{(n+1)(n+2)}\ge \dfrac{n}{6n^2}=\dfrac{1}{6n}.
\]
Hence the $N$th partial sum is at least $\dfrac{1}{6}\left(1+\dfrac{1}{2}+\cdots+\dfrac{1}{N}\right)$. The harmonic partial sums are unbounded: grouping the terms from $\dfrac{1}{2^{k-1}+1}$ to $\dfrac{1}{2^k}$ gives $2^{k-1}$ terms each at least $\dfrac{1}{2^k}$, so each group contributes at least $\dfrac{1}{2}$, and $1+\dfrac{1}{2}+\cdots+\dfrac{1}{2^k}\ge 1+\dfrac {k}{2}$.

Therefore the partial sums of the given series are unbounded, and
\[
\boxed{\text{the series diverges.}}
\]
